\section{Related Work}
\label{sec:related}

\paragraph{Testing role prompting assumptions.}
Studies show persona assignment is ineffectual at improving model accuracy, whereas \emph{oracle}
persona assignment is effective, suggesting the key is selection, not adoption
\citep{zheng2024helpful}. Similarly, studies using expert personas have mixed effects across
different tasks and outcomes \citep{hu2026expert}. Prompting models with different personalities
alters behaviour, with performance depending on the composition of the agent group
\citep{yu2025personality}. For this reason, there is reason for caution about role prompting,
rather than reasons to ensemble with it \citep{kim2024persona}. Any role prompting improvement
should be considered in this light.

\paragraph{Multi agent councils.}
Models are increasingly being ensembled across agents
\citep{bause2026mixtures,huang2025dynamics,hamidieh2026complementing}. The most direct comparison
is with \emph{Council Mode} \citep{wu2026council} which also shows reductions in hallucinations on
the same two datasets as we use. Their method combines three different models from three different
providers plus a specialised synthesising model; this means their improvement is a joint effect of
different specialisms, diversity of lineage, and synthesisation capability, which can't disentangle
each other. By contrast we keep the lineage constant and have no dedicated synthesiser,
differentiating each arm based solely on the prompt. We also have reason to be pessimistic about
council mode. The first is that deliberation may be an illusion \citep{wan2026deliberative}. Second
is that correlated failure limits improvement from combination
\citep{jamshidi2026cascade,chen2026cofailure}.

\paragraph{Self correction and verification.}
Self-correcting models have been shown to have questionable improvements, with self-correction
worsening answers \citep{huang2024selfcorrect}, self-verification failing to identify model
mistakes, and improvements being an illusion of evaluation
\citep{chen2026selfcorrection,stechly2024selfverification}. For this reason, we prefer to use
objective and actionable critiques over self critiques.

\paragraph{Testing for role adherence, closest prior work.}
Many existing works evaluate whether an LLM ``acts according to its assigned role'' on a prompt;
e.g.\ to what extent the LLM displays the \emph{persona} on psychological scales through interview
\citep{wang2024incharacter}; whether an LLM ``sticks to'' the \emph{persona} when completing
open-ended tasks \citep{shin2025spotting}; whether \emph{persona} parameters affect model
predictions at all \citep{hu2024persona}; specific checklists of role adherence for
\emph{role-playing agents} \citep{rosati2026rpacheck}; ``disobey role specification'' as a specific
failure mode for \emph{multi-agent systems} extracted from execution traces \citep{cemri2025why}.

We do not claim novelty, and there are two prior works sufficiently close to cite:
\citet{dietrich2026roles} construct role-specific tests of adherence to an advocate role using a
stance classifier, to ``explicitly test adherence to the respective roles;'' this is as close as we
can find to what we propose in a multi-agent LLM stack, and is a contemporaneous, non-peer-reviewed
preprint, and we weight it accordingly. \citet{bieberstein2026empathy} conduct an automatic
``manipulation check'' on a neutral prompt baseline \emph{before} reporting results in a
single-agent HCI study, which is identical to the procedure in psychology.

The novelty here does not lie in testing role adherence itself but in its application, in that the
aforementioned tests a) are not declared in advance; b) lack a baseline or length-unmatched
baseline; c) test adherence to a persona, not an explicitly-specified action; d) are never
interpreted as a precondition to interpreting a \textbf{benchmark}.

\paragraph{Architectural novelty, baselines, and provenance.}
Much prior work explores structured agent architectures \citep{sumers2024cognitive}, including
functional role agents governed by a routing agent \citep{shang2026theater} and role architectures
inspired by clinical frameworks \citep{wu2026wisemind} (see also \citet{lee2025emergence} for the
related research question of how and when agents ``spontaneously'' implement such computations). We
do not claim novelty on the agent design itself, but only on the quality control mechanism. Our
baseline uses cost-matched pre-registration, common for routing experiments \citep{ong2025routellm};
the results we obtain instead use a same-model single-call baseline (\S\ref{sec:arms}). The most
promising hope for functionally constraining LLMs lies in steering at the activation level
\citep{bas2025steer}, although such experiments have proved unreliable \citep{braun2026steering}; we
defer mechanistic intervention to future work. The archetypes we chose were initially inspired by
prior work on computational theories of mind and information processing-accounts in clinical
psychology and psychiatry
\citep{adams2016computational,addicott2021attention,clance1978impostor,frith2003autism,gruber2011bipolar,kapur2003psychosis}:
all such connections are limited to motivation of a functional role, and our archetypes are entirely
specified, implemented, and assessed in functional terms, and none of our conclusions apply to any
clinical group. \emph{These roles are functional AI processing constraints. They are not diagnostic
models, simulations of lived experience, or claims that mental health conditions confer a single
cognitive mechanism.}
